\providecommand{\SysName}{CARE}

\section{Introduction}
\label{sec:introduction}

Refactoring is a routine part of software maintenance: developers restructure
code to reduce duplication, simplify control flow, and improve long-term
maintainability without intending to change program behavior
\cite{fowler2018refactoring,mens2004survey}. In C/C++ systems code, however,
``behavior preserving'' is a stronger requirement than preserving ordinary
outputs. A small cleanup may change error codes, cleanup order, ownership
transfer, lock state, logging, or input-validation behavior. These details are
often security-relevant: a patch that removes a repeated check, rewrites an
error path, or consolidates duplicated cleanup can accidentally introduce a
resource leak, double release, unchecked allocation, weakened bounds check, or
changed API contract. Prior work has shown that refactoring can be useful but
risky in practice, and that testing and analysis are needed to make
refactorings safer \cite{soares2010making,kim2014empirical,vassallo2019large}.

Large language models (LLMs) create a new opportunity for automating such
refactorings. Modern code models can synthesize nontrivial code and are
increasingly studied for software transformation tasks
\cite{chen2021evaluating,li2022competition,liu2024empirical}. Yet direct LLM
patching is a poor fit for security-sensitive C/C++ maintenance. The model may
produce a plausible diff that does not apply, omit context from other control
paths, remove a check that is only locally redundant, or simplify cleanup in a
way that changes resource lifetime. The core problem is therefore not merely
how to ask an LLM to refactor code, but how to give it the right program
context, constrain the intended transformation, and reject unsafe candidates
before they reach maintainers.

This paper presents \SysName{}, a context-aware automated refactoring engine for
C/C++ code. \SysName{} takes a target project, builds a lightweight
security-aware context graph, detects refactoring opportunities, asks an LLM to
produce small behavior-preserving patch candidates, and validates each
candidate with a fail-closed pipeline. The current prototype targets
security-relevant refactorings common in systems code: duplicated cleanup
logic, resource lifecycle imbalance, repeated null/bounds/status checks, dead
or unreachable code, and general security smells. Rather than treating the LLM
as an autonomous developer, \SysName{} treats it as a proposal engine inside a
guarded synthesis loop: analysis supplies evidence and constraints, the LLM
generates candidate edits, and the validator decides whether the patch is safe
enough to report.

Designing such a system raises three challenges. First, \SysName{} must
construct enough program context to support refactoring decisions without
requiring full formal verification of an entire C/C++ project. It therefore
combines function extraction, lightweight control-flow construction, call-graph
edges, data-flow facts, resource events, error-handling paths, and side-effect
summaries into a security-aware contextual refactoring graph. Second,
\SysName{} must distinguish actionable opportunities from syntactic
coincidences. For example, a duplicate null check is safely removable only when
no assignment, mutating call, or obvious aliasing risk appears between the two
checks; similarly, a cleanup refactoring must preserve release order and return
conventions. Third, LLM-generated patches must be validated under the target
project's actual build and test environment. \SysName{} addresses this with
patch application checks, build/test execution, optional static analysis,
resource consistency checks, semantic-equivalence checks, security-regression
checks, and feedback-guided repair.

We implement \SysName{} as a modular Python research prototype. The system
supports project loading, clang-backed or fallback parsing, opportunity
detection, structured LLM planning, conservative and aggressive patch
generation modes, iterative repair, and report generation. We evaluate
\SysName{} on OSS50, a benchmark of 50 C/C++ open-source projects. A full scan
successfully covers 49 projects, 25,162 source files, and 266,602 functions,
producing 142,104 refactoring opportunities. For the completed critical
resource-lifecycle workload of 487 findings, \SysName{} improves patch
applicability by 3.1$\times$ and full validation yield by 6.2$\times$ over a
direct LLM-only baseline under the same LLM-call budget. The validator ablation
shows why fail-closed validation matters: removing the resource checker would
admit 158 additional patches, 121 of which are classified as false accepts by
review.

In summary, this paper makes the following contributions:
\begin{itemize}
    \item We introduce the problem of security-aware automated refactoring for
    C/C++ systems code, where behavior preservation must include resource
    lifecycles, error-handling behavior, validation checks, and externally
    visible side effects.
    \item We design \SysName{}, an end-to-end pipeline that combines program
    context construction, opportunity detection, LLM-based planning and patch
    generation, validation, feedback-guided repair, and reporting.
    \item We develop a security-aware contextual refactoring graph and detector
    suite for goto-based redundancy, resource imbalance, duplicate checks,
    dead/unreachable code, and security smells.
    \item We implement a fail-closed validation pipeline that checks patch
    application, build/test behavior, resource consistency, semantic stability,
    and security regressions before accepting an LLM-generated candidate.
    \item We evaluate \SysName{} on a large C/C++ open-source workload and show
    that context-aware generation plus validation substantially improves
    validated patch yield over direct LLM-only refactoring.
\end{itemize}
